% Chapter 3: System Analysis (Massive Expansion)
\chapter{SYSTEM ANALYSIS}

The success of any multi-cloud project depends on a rigorous analysis of the "Problem Space" before entering the "Solution Space." This chapter details the technical requirements, constraints, and the strategic planning involved in the Nebula platform.

\section{Software Requirements Specification (SRS)}
According to the IEEE-830 standard, we organized our requirements based on functional and non-functional goals.

\subsection{Actor Definitions}
1. \textbf{Cloud Administrator}: Responsible for adding cloud credentials and managing organizational access.
2. \textbf{DevOps Engineer}: The primary user who provisions, manages, and monitors resources across clouds.
3. \textbf{System Auditor}: Views cost reports and inventory history to ensure compliance.

\subsection{Use Case Descriptions}
Detailed use cases define how actors interact with the system to achieve specific business goals.

\begin{itemize}
    \item \textbf{Use Case 1: Secure Credential Registration}
        \begin{enumerate}
            \item User selects "Register Provider" and provides an Access Key and Secret Key.
            \item The system validates the format using Pydantic schemas.
            \item The backend encrypts keys using AES-128 and stores them in PostgreSQL.
            \item The user is redirected to the "Settings" dashboard.
        \end{enumerate}
    \item \textbf{Use Case 2: Multi-Cloud VM Provisioning}
        \begin{enumerate}
            \item User clicks "New VM" and selects "AWS" as the provider.
            \item System fetches the predefined "EC2" Terraform module.
            \item User enters instance name, size, and region.
            \item System initiates a Celery task and provides a real-time status update.
        \end{enumerate}
\end{itemize}

\section{Requirement Traceability Matrix (RTM)}
The RTM is a document that maps and traces user requirement with test cases. It captures all requirements proposed by the client and requirement traceability in a single document, delivered at the conclusion of the Software development life cycle. 

\begin{table}[H]
\centering
\caption{Requirement Traceability Matrix}
\label{tab:rtm}
\begin{tabular}{|p{1.5cm}|p{4.5cm}|p{3.5cm}|p{2cm}|}
\hline
\textbf{Req ID} & \textbf{Requirement Description} & \textbf{Design Ref} & \textbf{Test ID} \\ \hline
FR-01 & User Authentication via JWT & Sec-Arch 4.5 & TC-01 \\ \hline
FR-02 & Encrypted Credential Vault & Sec-Arch 4.5 & TC-02 \\ \hline
FR-03 & AWS Resource Provisioning & Engine 5.3 & TC-03 \\ \hline
FR-04 & Azure Resource Provisioning & Engine 5.3 & TC-04 \\ \hline
FR-05 & GCP Resource Provisioning & Engine 5.3 & TC-05 \\ \hline
FR-06 & AI Log Remediation Engine & AI-Logic 8.1 & TC-06 \\ \hline
FR-07 & Real-time Status Sync & Sync-Arch 4.6 & TC-07 \\ \hline
NFR-01 & Encryption at Rest (AES-128) & Sec-Arch 4.5 & TC-14 \\ \hline
NFR-02 & API Latency < 1s & HLD 4.1 & TC-33 \\ \hline
NFR-03 & System Uptime 99.9\% & Ops-Manual & TC-35 \\ \hline
\end{tabular}
\end{table}

\section{Detailed Project Schedule (16-Week Plan)}
The project was executed following a structured timeline, ensuring that every phase of the SDLC was given adequate time for research, design, implementation, and testing.

\subsection{Phase 1: Research and Feasibility (Weeks 1-3)}
- Identification of multi-cloud management gaps.
- Survey of 10+ cloud tools (as detailed in Chapter 2).
- Validation of Terraform as the core orchestration engine.

\subsection{Phase 2: System Architecture Design (Weeks 4-6)}
- Formalization of the High-Level Design (HLD).
- Selection of tech stack: FastAPI, React, PostgreSQL, Redis, and Celery.
- Database schema design and entity definition.

\subsection{Phase 3: Core Module Implementation (Weeks 7-11)}
- Development of the "Nebula Vault" for credential security.
- Implementation of the Terraform Task Runner.
- Creation of provider-specific HCL modules (AWS, Azure, GCP).

\subsection{Phase 4: Frontend Development (Weeks 12-14)}
- Building the "Command Center" Dashboard.
- Integration with backend WebSockets for real-time log streaming.
- Implementation of the AI-powered Copilot UI.

\subsection{Phase 5: Testing and Deployment (Weeks 15-16)}
- Functional and Non-Functional testing (as detailed in Chapter 6).
- Security auditing of the credential encryption logic.
- Final documentation and user manual preparation.

\section{Risk Management and Mitigation}
Managing multiple cloud providers introduces unique risks that were identified and addressed during the analysis phase.

\begin{table}[H]
\centering
\caption{Project Risk Matrix}
\label{tab:risk_matrix}
\begin{tabular}{|p{3cm}|p{1.5cm}|p{4.5cm}|p{3.5cm}|}
\hline
\textbf{Risk} & \textbf{Impact} & \textbf{Mitigation Strategy} & \textbf{Contingency} \\ \hline
Cloud API Rate Limiting & High & Implement exponential backoff in sync tasks. & Alert administrator. \\ \hline
Credential Compromise & Critical & Encryption at application layer (AES-128). & Secret rotation policy. \\ \hline
Terraform State Corruption & High & Use local backups and isolation per resource. & State recovery tools. \\ \hline
AI Hallucination & Medium & Sanitized logs and human-in-the-loop validation. & Fallback to raw logs. \\ \hline
Provider Outage & High & Multi-cloud redundancy (Nebula survives if 1 cloud fails). & Failover to other cloud. \\ \hline
\end{tabular}
\end{table}

\section{Data Dictionary for Core Entities}
The following data dictionary describes the attributes of each operational entity in the PostgreSQL schema.

\begin{table}[H]
\centering
\caption{Nebula Data Dictionary}
\label{tab:dictionary}
\begin{tabular}{|l|l|l|l|}
\hline
\textbf{Entity} & \textbf{Attribute} & \textbf{Type} & \textbf{Description} \\ \hline
Credential & provider\_type & Enum & AWS, AZURE, or GCP \\ \hline
Credential & access\_key & String (Enc) & Encrypted primary access key \\ \hline
Resource & provider\_id & String & Unique cloud ID (e.g., i-12345) \\ \hline
Resource & status & Enum & active, pending, failed, or deleted \\ \hline
Task & execution\_log & Text & Full STDOUT/STDERR from Terraform \\ \hline
\end{tabular}
\end{table}

\newpage
